\begin{afterword}
\summaryhead

The post quotes two comments by Kurige. In the first, Kurige says that a world-view once based on blind faith has moved to ``somewhere in the vicinity of Orwellian double-think.'' In the second, Kurige says of believing in God: ``I chose to believe \ldots\ deliberately and consciously,'' and the decision ``has absolutely zero effect on the actual existence of God.'' The author asks how anyone can believe what they know is double-think, or a belief they know is not correlated with reality.

Part of the answer may be the account given in ``Moore's Paradox.'' Another part may be that it is ``quite easy'' not to go from knowing what an accurate map would say to believing it. The author now says that the earlier sentence ``You cannot make yourself believe the sky is green by an act of will'' was partly an attempt to instill a self-fulfilling prophecy. The post recommends repeating to yourself that deliberate self-deception cannot work, since believing in your own inability makes you less likely to succeed. Positive or negative thinking, it seems a wise precaution.

\responsehead

The most important sentence in the post is a partial retraction, and it is made openly. ``Doublethink'' (2007) offered ``You cannot make yourself believe the sky is green by an act of will'' as fact, and ``Belief in Self-Deception'' repeated it four days before this post. Here the author says the sentence was not only a report of ``the existing facts'' but also an attempt to make itself true, and grants that it ``may just be'' quite easy not to believe what the accurate map says. In the book, readers meet the confident version in ``Doublethink'' and this admission four posts later.

The puzzle the post starts from rests on a reading of Kurige's comments that they do not support. Kurige used ``double-think'' loosely, for holding two accounts of morality, and said in the next paragraph that ``These two beliefs are not contradictory.'' And the decision's ``zero effect on the actual existence of God'' is true of every belief, accurate ones included: on the author's own account, a belief tracks reality by being caused by it, not by causing it. The same comment compared the belief to a scientist's belief in the Higgs boson, held on evidence. So the post's question, how one can believe what one knows is uncorrelated with reality, is not one Kurige's words raise.

The advice has the form the earlier posts argued against: adopt a belief for its effects. ``Doublethink'' opposed the idea that you can ``decide which cognitive biases should govern you.'' What differs here is that the belief is meant to make itself true once held. That defence depends on the effect being real, and the post asserts it without a hedge or evidence: ``you will indeed be less likely to fool yourself successfully.'' The hedges sit on the recommendation instead (``It may be wise,'' ``it seems like a wise precaution''), and the slogan drops them: ``Tell yourself the effort is doomed---and it will be!''

In short: an open admission that an earlier confident claim was partly a self-fulfilling prophecy, a puzzle built on a reading of Kurige's comments that they do not support, and advice to hold a belief for its effects, with the effect asserted rather than shown.
\end{afterword}
